% !TeX program = xelatex
% UTF-8; compile twice with XeLaTeX. No external images or data files needed.
\documentclass[a4paper,landscape,10pt]{article}
\usepackage[margin=14mm,headheight=14pt,headsep=5mm,footskip=7mm]{geometry}
\usepackage{fontspec}
\setmainfont{DejaVu Sans}
\newfontfamily\ipafont{DejaVu Sans}
\usepackage[table]{xcolor}
\usepackage{array,longtable,booktabs,fancyhdr}
\definecolor{ink}{HTML}{18384D}
\definecolor{band}{HTML}{E7EFF4}
\pagestyle{fancy}\fancyhf{}
\fancyhead[L]{\small\color{ink}Parlez Vous QSO?! · Українська / Ukrainian}
\fancyhead[R]{\small Englische Phrase · Ukrainische Übersetzung · IPA}
\fancyfoot[L]{\scriptsize Neufassung: Ukrainisch · Projektidee und historische Dokumente: Steve -- M0ECS}
\fancyfoot[R]{\small\thepage}
\setlength{\parindent}{0pt}
\setlength{\tabcolsep}{3mm}
\renewcommand{\arraystretch}{1.25}
\newcommand{\ipa}[1]{{\ipafont[#1]}}
\begin{document}
{\LARGE\bfseries\color{ink}Parlez Vous QSO?!}\hfill{\large Українська / Ukrainian}

\vspace{2mm}
{\small QSO-Sprachführer mit ukrainischer IPA-Lautschrift. Die Lautschrift bezieht sich auf die ukrainische Spalte.}

\vspace{2mm}
{\footnotesize \textbf{Leseschlüssel:} \ipa{ˈ} Hauptbetonung; \ipa{ʲ} bezeichnet einen weichen (palatalisierten) Konsonanten.
\ipa{ɦ} ist ein stimmhafter h-Laut, \ipa{ɡ} wie deutsches g; \ipa{x} wie „ch“ in „Bach“, \ipa{ʃ} wie „sch“.
\ipa{ɪ} entspricht ukrainischem и, \ipa{i} ukrainischem і. Grundlage: ukrainische Standardsprache.}

\vspace{1mm}
\fontsize{9}{11.8}\selectfont
\begin{longtable}{@{}>{\raggedright\arraybackslash}p{13mm}>{\raggedright\arraybackslash}p{73mm}>{\raggedright\arraybackslash}p{75mm}>{\raggedright\arraybackslash}p{90mm}@{}}
\toprule
\textbf{Ref.} & \textbf{English phrase} & \textbf{Українська} & \textbf{Lautschrift (IPA)}\\
\midrule
\endfirsthead
\toprule
\textbf{Ref.} & \textbf{English phrase} & \textbf{Українська} & \textbf{Lautschrift (IPA)}\\
\midrule
\endhead
\midrule
\multicolumn{4}{r}{\scriptsize Fortsetzung auf der nächsten Seite}\\
\endfoot
\bottomrule
\endlastfoot
\rowcolor{band} 1.0 & \textbf{Opening Greetings} & \multicolumn{2}{l}{\textbf{Початкові привітання}} \\*
1.1 & I am very pleased to meet you. & Дуже радий нашому знайомству. & \ipa{ˈduʒɛ ˈradɪj ˈnaʃɔmu znaˈjɔmstʋu} \\
1.2 & Thank you for your call & Дякую за Ваш виклик. & \ipa{ˈdʲakuju za ʋaʃ ˈʋɪklɪk} \\
1.3 & Good morning & Доброго ранку. & \ipa{ˈdɔbrɔɦɔ ˈranku} \\
1.31 & Good day & Добрий день. & \ipa{ˈdɔbrɪj dɛnʲ} \\
1.32 & Good afternoon & Добрий день. & \ipa{ˈdɔbrɪj dɛnʲ} \\
1.33 & Good evening & Добрий вечір. & \ipa{ˈdɔbrɪj ˈʋɛtʃʲir} \\
1.4 & I know only a little Ukrainian but I enjoy QSOs in Ukrainian. & Я знаю українську лише трохи, але люблю проводити QSO українською. & \ipa{ja ˈznaju ukraˈjinʲsʲku lɪˈʃɛ ˈtrɔxɪ aˈlɛ lʲubˈlʲu prɔˈʋɔdɪtɪ kjuː ɛs oʊ ukraˈjinʲsʲkɔju} \\
\rowcolor{band} 2.0 & \textbf{Operator} & \multicolumn{2}{l}{\textbf{Оператор}} \\*
2.1 & My name is Peter. & Мене звати Пітер. & \ipa{mɛˈnɛ ˈzʋatɪ ˈpʲitɛr} \\
2.2 & I have been licensed since 1960. & Маю ліцензію з 1960 року. & \ipa{ˈmaju lʲiˈt͡sɛnzʲiju z ˈtɪsʲatʃa dɛʋjatʲˈsɔt ʃizdɛˈsʲatɔɦɔ ˈrɔku} \\
2.3 & I have been licensed for 43 years. & Маю ліцензію вже 43 роки. & \ipa{ˈmaju lʲiˈt͡sɛnzʲiju ʋʒɛ ˈsɔrɔk trɪ ˈrɔkɪ} \\
2.4 & I live in Sittingbourne. & Я живу в Сіттінгборні. & \ipa{ja ʒɪˈʋu ʋ ˈsʲitʲinɦbɔrnʲi} \\
2.5 & Sittingbourne is a village. & Сіттінгборн - це село. & \ipa{ˈsʲitʲinɦbɔrn t͡sɛ sɛˈlɔ} \\
2.51 & Sittingbourne is a town. & Сіттінгборн - це містечко. & \ipa{ˈsʲitʲinɦbɔrn t͡sɛ mʲiˈstɛtʃkɔ} \\
2.52 & Sittingbourne is a city. & Сіттінгборн - це місто. & \ipa{ˈsʲitʲinɦbɔrn t͡sɛ ˈmʲistɔ} \\
2.6 & Sittingbourne is between London and Dover. & Сіттінгборн розташований між Лондоном і Дувром. & \ipa{ˈsʲitʲinɦbɔrn rɔztaˈʃɔʋanɪj mʲiʒ ˈlɔndɔnɔm i ˈduʋrɔm} \\
2.7 & I live in the South of England. & Я живу на півдні Англії. & \ipa{ja ʒɪˈʋu na pʲiʋˈdnʲi ˈanɦlʲiji} \\
2.71 & I live in the North of England. & Я живу на півночі Англії. & \ipa{ja ʒɪˈʋu na ˈpʲiʋnɔtʃʲi ˈanɦlʲiji} \\
2.72 & I live in the East of England. & Я живу на сході Англії. & \ipa{ja ʒɪˈʋu na ˈsxɔdʲi ˈanɦlʲiji} \\
2.73 & I live in West of England. & Я живу на заході Англії. & \ipa{ja ʒɪˈʋu na ˈzaxɔdʲi ˈanɦlʲiji} \\
2.74 & I live in South East England. & Я живу на південному сході Англії. & \ipa{ja ʒɪˈʋu na pʲiʋˈdɛnnɔmu ˈsxɔdʲi ˈanɦlʲiji} \\
2.75 & I live in South West England. & Я живу на південному заході Англії. & \ipa{ja ʒɪˈʋu na pʲiʋˈdɛnnɔmu ˈzaxɔdʲi ˈanɦlʲiji} \\
2.76 & I live in North East England. & Я живу на північному сході Англії. & \ipa{ja ʒɪˈʋu na pʲiʋˈnʲitʃnɔmu ˈsxɔdʲi ˈanɦlʲiji} \\
2.77 & I live in North West England. & Я живу на північному заході Англії. & \ipa{ja ʒɪˈʋu na pʲiʋˈnʲitʃnɔmu ˈzaxɔdʲi ˈanɦlʲiji} \\
2.78 & I live in central England. & Я живу в центральній Англії. & \ipa{ja ʒɪˈʋu ʋ t͡sɛnˈtralʲnʲij ˈanɦlʲiji} \\
2.8 & I live in the country. & Я живу в сільській місцевості. & \ipa{ja ʒɪˈʋu ʋ ˈsʲilʲsʲkʲij mʲisˈt͡sɛʋɔstʲi} \\
2.81 & I live by the sea. & Я живу біля моря. & \ipa{ja ʒɪˈʋu ˈbʲilʲa ˈmɔrʲa} \\
2.82 & I live near mountains. & Я живу біля гір. & \ipa{ja ʒɪˈʋu ˈbʲilʲa ɦʲir} \\
2.9 & Sittingbourne is 50Km from London. & Сіттінгборн розташований за 50 км від Лондона. & \ipa{ˈsʲitʲinɦbɔrn rɔztaˈʃɔʋanɪj za pjatdɛˈsʲat kʲilɔˈmɛtrʲiʋ ʋʲid ˈlɔndɔna} \\
2.11 & I am retired. & Я на пенсії. & \ipa{ja na ˈpɛnsʲiji} \\
2.12 & My home call sign is M0ECS & Мій домашній позивний - M0ECS. & \ipa{mʲij dɔˈmaʃnʲij pɔzɪʋˈnɪj ɛm nulʲ iː siː ɛs} \\
\rowcolor{band} 3.0 & \textbf{Station} & \multicolumn{2}{l}{\textbf{Радіостанція}} \\*
3.1 & My radio is an Icom 706. & Мій трансивер - Icom 706. & \ipa{mʲij tranˈsɪʋɛr ˈaɪkɒm sʲimˈsɔt ʃʲistʲ} \\
3.2 & My station is all homebrew (homemade). & Уся моя радіостанція саморобна. & \ipa{uˈsʲa mɔˈja radʲiɔˈstant͡sʲija samɔˈrɔbna} \\
3.3 & I am running 5 Watts. & Я працюю з потужністю 5 ват. & \ipa{ja praˈt͡sʲuju z pɔˈtuʒnʲistʲu pjatʲ ʋat} \\
3.4 & My antenna is a long wire. & Моя антена - довгий дріт. & \ipa{mɔˈja anˈtɛna ˈdɔʋɦɪj drʲit} \\
3.41 & My antenna is a dipole. & Моя антена - диполь. & \ipa{mɔˈja anˈtɛna dɪˈpɔlʲ} \\
3.42 & My antenna is a G5RV. & Моя антена - G5RV. & \ipa{mɔˈja anˈtɛna dʒiː pjatʲ ɑː viː} \\
3.43 & My antenna is a doublet. & Моя антена - диполь типу doublet. & \ipa{mɔˈja anˈtɛna dɪˈpɔlʲ ˈtɪpu ˈdʌblɪt} \\
3.44 & My antenna is a delta loop. & Моя антена - трикутна рамка. & \ipa{mɔˈja anˈtɛna trɪˈkutna ˈramka} \\
3.45 & My antenna is a vertical. & Моя антена вертикальна. & \ipa{mɔˈja anˈtɛna ʋɛrtɪˈkalʲna} \\
3.46 & My antenna is a log periodic. & Моя антена логоперіодична. & \ipa{mɔˈja anˈtɛna lɔɦɔpɛrʲiɔˈdɪtʃna} \\
3.47 & My antenna is a Yagi beam. & Моя антена - спрямована Яґі. & \ipa{mɔˈja anˈtɛna sprʲaˈmɔʋana ˈjaɡʲi} \\
3.48 & My antenna is a 5 element beam. & Моя антена спрямована, з 5 елементами. & \ipa{mɔˈja anˈtɛna sprʲaˈmɔʋana z pjatʲˈma ɛlɛˈmɛntamɪ} \\
\rowcolor{band} 4.0 & \textbf{Conditions} & \multicolumn{2}{l}{\textbf{Умови поширення радіохвиль}} \\*
4.1 & Thank you for my 599 report. & Дякую за рапорт 599. & \ipa{ˈdʲakuju za ˈrapɔrt pjatʲ ˈdɛʋjatʲ ˈdɛʋjatʲ} \\
4.2 & The band conditions are good. & Умови проходження на діапазоні хороші. & \ipa{uˈmɔʋɪ prɔˈxɔdʒɛnnʲa na dʲiapaˈzɔnʲi xɔˈrɔʃʲi} \\
4.3 & The band conditions are bad. & Умови проходження на діапазоні погані. & \ipa{uˈmɔʋɪ prɔˈxɔdʒɛnnʲa na dʲiapaˈzɔnʲi pɔˈɦanʲi} \\
4.4 & There is strong interference here. & Тут сильні завади. & \ipa{tut ˈsɪlʲnʲi zaˈʋadɪ} \\
4.5 & Can you change to the 80m band? & Чи можете перейти на діапазон 80 метрів? & \ipa{tʃɪ ˈmɔʒɛtɛ pɛrɛjˈtɪ na dʲiapaˈzɔn ʋʲisʲimdɛˈsʲat ˈmɛtrʲiʋ} \\
4.6 & Please change frequency to 3.558 MHz. & Будь ласка, перейдіть на частоту 3,558 МГц. & \ipa{budʲ ˈlaska pɛrɛjˈdʲitʲ na tʃastɔˈtu trɪ ˈkɔma pjatʲ pjatʲ ˈʋʲisʲim mɛɦaˈɦɛrt͡s} \\
\rowcolor{band} 5.0 & \textbf{General Queries and Comments} & \multicolumn{2}{l}{\textbf{Загальні запитання та коментарі}} \\*
5.1 & What does <\#\#\#\#\#> mean please? & Що означає <\#\#\#\#\#>, будь ласка? & \ipa{ʃtʃɔ ɔznaˈtʃajɛ … budʲ ˈlaska} \\
5.2 & Please wait a moment, I need to check my Ukrainian dictionary. & Зачекайте, будь ласка. Мені треба заглянути в український словник. & \ipa{zatʃɛˈkajtɛ budʲ ˈlaska mɛˈnʲi ˈtrɛba zaˈɦlʲanutɪ ʋ ukraˈjinʲsʲkɪj slɔʋˈnɪk} \\
5.3 & How do you say, <\#\#\#\#\#> in Ukrainian ? & Як сказати <\#\#\#\#\#> українською? & \ipa{jak skaˈzatɪ … ukraˈjinʲsʲkɔju} \\
5.5 & My favourite modes of operation are CW / RTTY / PSK / MT63 / Olivia / Hellschreiber / Digital Modes / and SSB & Мої улюблені види роботи - CW / RTTY / PSK / MT63 / Olivia / Hellschreiber / цифрові види / та SSB. & \ipa{mɔˈji uˈlʲublɛnʲi ˈʋɪdɪ rɔˈbɔtɪ siː ˈdʌbljuː / ɑː tiː tiː waɪ / piː ɛs keɪ / ɛm tiː ʃizdɛˈsʲat trɪ / ɔˈlʲiʋʲija / ˈhɛlˌʃʁaɪbɐ / t͡sɪfrɔˈʋʲi ˈʋɪdɪ / ta ɛs ɛs biː} \\
5.6 & My best DX is Australia. & Мій найкращий DX - Австралія. & \ipa{mʲij najˈkraʃtʃɪj diː ɛks aʋˈstralʲija} \\
5.7 & I enjoy working contests. & Люблю брати участь у змаганнях. & \ipa{lʲubˈlʲu ˈbratɪ ˈutʃastʲ u zmaˈɦannʲax} \\
5.8 & I like working QRP. & Люблю працювати в режимі QRP. & \ipa{lʲubˈlʲu prat͡sʲuˈʋatɪ ʋ rɛˈʒɪmʲi kjuː ɑː piː} \\
\rowcolor{band} 6.0 & \textbf{The Weather} & \multicolumn{2}{l}{\textbf{Погода}} \\*
6.1 & It is sunny. & Сонячно. & \ipa{ˈsɔnʲatʃnɔ} \\
6.2 & It is windy. & Вітряно. & \ipa{ˈʋʲitrʲanɔ} \\
6.3 & It is raining. & Іде дощ. & \ipa{iˈdɛ dɔʃtʃ} \\
6.4 & It is snowing. & Іде сніг. & \ipa{iˈdɛ sʲnʲiɦ} \\
6.5 & It is very hot. & Дуже спекотно. & \ipa{ˈduʒɛ spɛˈkɔtnɔ} \\
6.6 & It is very cold. & Дуже холодно. & \ipa{ˈduʒɛ ˈxɔlɔdnɔ} \\
6.7 & The temperature is 8C. & Температура - 8 °C. & \ipa{tɛmpɛraˈtura ˈʋʲisʲim ˈɦradusʲiʋ ˈt͡sɛlʲsʲija} \\
6.8 & It is overcast. & Небо вкрите хмарами. & \ipa{ˈnɛbɔ ˈʋkrɪtɛ ˈxmaramɪ} \\
6.9 & It is foggy. & Туманно. & \ipa{tuˈmannɔ} \\
6.11 & It has been raining all day. & Дощ іде вже весь день. & \ipa{dɔʃtʃ iˈdɛ ʋʒɛ ʋɛsʲ dɛnʲ} \\
6.12 & There is 20cm of snow here. & Тут 20 см снігу. & \ipa{tut ˈdʋad͡zʲt͡sʲatʲ santɪˈmɛtrʲiʋ ˈsʲnʲiɦu} \\
\rowcolor{band} 8.0 & \textbf{QSL} & \multicolumn{2}{l}{\textbf{Картки QSL}} \\*
8.1 & My QSL card is sure via the bureau & Обов'язково надішлю свою картку QSL через бюро. & \ipa{ɔbɔʋjazˈkɔʋɔ nadʲiˈʃlʲu sʋɔˈju ˈkartku kjuː ɛs ɛl ˈtʃɛrɛz bʲuˈrɔ} \\
8.2 & QSL via eQSL & QSL через eQSL. & \ipa{kjuː ɛs ɛl ˈtʃɛrɛz iː kjuː ɛs ɛl} \\
8.3 & No eQSL please & Будь ласка, не надсилайте eQSL. & \ipa{budʲ ˈlaska nɛ nadsɪˈlajtɛ iː kjuː ɛs ɛl} \\
8.4 & I will send you my card direct. & Надішлю Вам свою картку безпосередньо. & \ipa{nadʲiˈʃlʲu ʋam sʋɔˈju ˈkartku bɛzpɔsɛˈrɛdʲnʲɔ} \\
8.5 & What is your address? & Яка Ваша адреса? & \ipa{jaˈka ˈʋaʃa aˈdrɛsa} \\
8.6 & Please send your QSL card direct. & Будь ласка, надішліть свою картку QSL безпосередньо. & \ipa{budʲ ˈlaska nadʲiʃˈlʲitʲ sʋɔˈju ˈkartku kjuː ɛs ɛl bɛzpɔsɛˈrɛdʲnʲɔ} \\
8.7 & My address is... & Моя адреса... & \ipa{mɔˈja aˈdrɛsa …} \\
8.8 & My address is shown on QSL.com / EHam.net & Моя адреса є на QSL.com / EHam.net. & \ipa{mɔˈja aˈdrɛsa jɛ na kjuː ɛs ɛl ˈkrapka kɔm / iː hæm ˈkrapka nɛt} \\
8.9 & Sorry I do not send QSL cards. & Вибачте, я не надсилаю карток QSL. & \ipa{ˈʋɪbatʃtɛ ja nɛ nadsɪˈlaju karˈtɔk kjuː ɛs ɛl} \\
8.11 & Please send QSL card via G1JQH. & Будь ласка, надішліть картку QSL через G1JQH. & \ipa{budʲ ˈlaska nadʲiʃˈlʲitʲ ˈkartku kjuː ɛs ɛl ˈtʃɛrɛz dʒiː ɔˈdɪn dʒeɪ kjuː eɪtʃ} \\
8.12 & My WACRAL / GQRP / DXCC / FISTS number is 1234. & Мій номер WACRAL / GQRP / DXCC / FISTS - 1234. & \ipa{mʲij ˈnɔmɛr ˈdʌbljuː eɪ siː ɑː eɪ ɛl / dʒiː kjuː ɑː piː / diː ɛks siː siː / fɪsts ˈtɪsʲatʃa ˈdʋistʲi ˈtrɪd͡zʲt͡sʲatʲ tʃɔˈtɪrɪ} \\
\rowcolor{band} 9.0 & \textbf{Final Greetings} & \multicolumn{2}{l}{\textbf{Завершальні привітання}} \\*
9.1 & Thank you. & Дякую. & \ipa{ˈdʲakuju} \\
9.2 & Thank you for a very enjoyable contact. & Дякую за дуже приємний зв'язок. & \ipa{ˈdʲakuju za ˈduʒɛ prɪˈjɛmnɪj zʋjaˈzɔk} \\
9.3 & I enjoyed our chat. & Мені сподобалася наша розмова. & \ipa{mɛˈnʲi spɔˈdɔbalasʲa ˈnaʃa rɔzˈmɔʋa} \\
9.4 & I enjoyed our QSO in Ukrainian. & Мені сподобалося наше QSO українською. & \ipa{mɛˈnʲi spɔˈdɔbalɔsʲa ˈnaʃɛ kjuː ɛs oʊ ukraˈjinʲsʲkɔju} \\
9.5 & It has been great fun. & Було дуже весело. & \ipa{buˈlɔ ˈduʒɛ ˈʋɛsɛlɔ} \\
9.6 & If you would like a regular contact with me, you can contact me via email. & Якщо хочете підтримувати регулярний зв'язок зі мною, можете писати мені електронною поштою. & \ipa{jakˈʃtʃɔ ˈxɔtʃɛtɛ pʲidˈtrɪmuʋatɪ rɛɦuˈlʲarnɪj zʋjaˈzɔk zʲi ˈmnɔju ˈmɔʒɛtɛ pɪˈsatɪ mɛˈnʲi ɛlɛkˈtrɔnnɔju ˈpɔʃtɔju} \\
9.7 & My email address is... & Моя електронна адреса... & \ipa{mɔˈja ɛlɛkˈtrɔnna aˈdrɛsa …} \\
9.8 & My email address is shown on QSL.com / EHam.net & Моя електронна адреса є на QSL.com / EHam.net. & \ipa{mɔˈja ɛlɛkˈtrɔnna aˈdrɛsa jɛ na kjuː ɛs ɛl ˈkrapka kɔm / iː hæm ˈkrapka nɛt} \\
9.9 & I look forward to meeting you again. & Сподіваюся знову зустрітися з Вами. & \ipa{spɔdʲiˈʋajusʲa ˈznɔʋu zuˈstrʲitɪsʲa z ˈʋamɪ} \\
9.11 & My best wishes to you and your family. & Найкращі побажання Вам і Вашій родині. & \ipa{najˈkraʃtʃʲi pɔbaˈʒannʲa ʋam i ˈʋaʃʲij rɔˈdɪnʲi} \\
9.12 & Best wishes to your wife. & Найкращі побажання Вашій дружині. & \ipa{najˈkraʃtʃʲi pɔbaˈʒannʲa ˈʋaʃʲij druˈʒɪnʲi} \\
9.121 & Best wishes to your husband. & Найкращі побажання Вашому чоловікові. & \ipa{najˈkraʃtʃʲi pɔbaˈʒannʲa ˈʋaʃɔmu tʃɔlɔˈʋʲikɔʋʲi} \\
9.13 & Good luck. & Успіхів. & \ipa{uˈspʲixʲiʋ} \\
9.14 & Good health. & Міцного здоров'я. & \ipa{mʲit͡sˈnɔɦɔ zdɔˈrɔʋja} \\
9.15 & Goodbye & До побачення. & \ipa{dɔ pɔˈbatʃɛnnʲa} \\
9.16 & I wish you an enjoyable afternoon. & Бажаю Вам гарного дня. & \ipa{baˈʒaju ʋam ˈɦarnɔɦɔ dnʲa} \\
9.161 & I wish you an enjoyable evening. & Бажаю Вам гарного вечора. & \ipa{baˈʒaju ʋam ˈɦarnɔɦɔ ˈʋɛtʃɔra} \\
9.162 & I wish you an enjoyable weekend. & Бажаю Вам гарних вихідних. & \ipa{baˈʒaju ʋam ˈɦarnɪx ʋɪxʲidˈnɪx} \\
9.17 & Enjoy your holiday. & Гарної відпустки. & \ipa{ˈɦarnɔji ʋʲidˈpustkɪ} \\
9.18 & Happy New Year. & З Новим роком. & \ipa{z nɔˈʋɪm ˈrɔkɔm} \\
\rowcolor{band} 10.0 & \textbf{Country} & \multicolumn{2}{l}{\textbf{Країни та території}} \\*
10.1 & Albania & Албанія & \ipa{alˈbanʲija} \\
10.2 & Austria & Австрія & \ipa{ˈaʋstrʲija} \\
10.3 & Belgium & Бельгія & \ipa{ˈbɛlʲɦʲija} \\
10.4 & Belorussia & Білорусь & \ipa{bʲilɔˈrusʲ} \\
10.5 & Bosnia \& Herzegovina & Боснія і Герцеговина & \ipa{ˈbɔsnʲija i ɦɛrt͡sɛɦɔˈʋɪna} \\
10.6 & Bulgaria & Болгарія & \ipa{bɔlˈɦarʲija} \\
10.7 & Croatia & Хорватія & \ipa{xɔrˈʋatʲija} \\
10.8 & Czech Republic & Чехія & \ipa{ˈtʃɛxʲija} \\
10.9 & Denmark & Данія & \ipa{ˈdanʲija} \\
10.11 & Eire & Ірландія & \ipa{irˈlandʲija} \\
10.12 & England & Англія & \ipa{ˈanɦlʲija} \\
10.13 & Estonia & Естонія & \ipa{ɛsˈtɔnʲija} \\
10.14 & Finland & Фінляндія & \ipa{fʲinˈlʲandʲija} \\
10.15 & France & Франція & \ipa{ˈfrant͡sʲija} \\
10.16 & Germany & Німеччина & \ipa{nʲiˈmɛtʃtʃɪna} \\
10.17 & Gibraltar & Гібралтар & \ipa{ɦʲibralˈtar} \\
10.18 & Greece & Греція & \ipa{ˈɦrɛt͡sʲija} \\
10.19 & Holland & Голландія & \ipa{ɦɔlˈlandʲija} \\
10.21 & Hungary & Угорщина & \ipa{ˈuɦɔrʃtʃɪna} \\
10.22 & Iceland & Ісландія & \ipa{isˈlandʲija} \\
10.23 & Northern Ireland & Північна Ірландія & \ipa{pʲiʋˈnʲitʃna irˈlandʲija} \\
10.24 & Isle of Man & Острів Мен & \ipa{ˈɔstrʲiʋ mɛn} \\
10.25 & Italy & Італія & \ipa{iˈtalʲija} \\
10.26 & Latvia & Латвія & \ipa{ˈlatʋʲija} \\
10.27 & Liechtenstein & Ліхтенштейн & \ipa{lʲixtɛnˈʃtɛjn} \\
10.28 & Lithuania & Литва & \ipa{lɪtˈʋa} \\
10.29 & Luxembourg & Люксембург & \ipa{lʲuksɛmˈburɦ} \\
10.31 & Macedonia & Північна Македонія & \ipa{pʲiʋˈnʲitʃna makɛˈdɔnʲija} \\
10.32 & Moldavia & Молдова & \ipa{mɔlˈdɔʋa} \\
10.33 & Netherlands & Нідерланди & \ipa{nʲidɛrˈlandɪ} \\
10.34 & Norway & Норвегія & \ipa{nɔrˈʋɛɦʲija} \\
10.35 & Poland & Польща & \ipa{ˈpɔlʲʃtʃa} \\
10.36 & Portugal & Португалія & \ipa{pɔrtuˈɦalʲija} \\
10.37 & Romania & Румунія & \ipa{ruˈmunʲija} \\
10.38 & Russian Federation & Російська Федерація & \ipa{rɔˈsʲijsʲka fɛdɛˈrat͡sʲija} \\
10.39 & Scotland & Шотландія & \ipa{ʃɔtˈlandʲija} \\
10.41 & Slovakia & Словаччина & \ipa{slɔˈʋatʃtʃɪna} \\
10.42 & Slovenia & Словенія & \ipa{slɔˈʋɛnʲija} \\
10.43 & Spain & Іспанія & \ipa{isˈpanʲija} \\
10.44 & Svalbard & Шпіцберген & \ipa{ʃpʲit͡sˈbɛrɦɛn} \\
10.45 & Sweden & Швеція & \ipa{ˈʃʋɛt͡sʲija} \\
10.46 & Switzerland & Швейцарія & \ipa{ʃʋɛjˈt͡sarʲija} \\
10.47 & Turkey & Туреччина & \ipa{tuˈrɛtʃtʃɪna} \\
10.48 & Ukraine & Україна & \ipa{ukraˈjina} \\
10.49 & Wales & Уельс & \ipa{uˈɛlʲs} \\
\end{longtable}
\clearpage
\normalsize
\section*{Hinweise zur Neufassung und Aussprache}
\fontsize{8.5}{10.5}\selectfont\raggedright
\textbf{Herkunft und Umfang.} Die ursprüngliche Idee zu „Parlez Vous QSO“ und die wesentlichen
historischen DOC-Dokumente stammen von Steve -- M0ECS. Michael Renner hat das Projekt mit Steves
Einverständnis übernommen und führt es fort. Diese ukrainische Fassung wurde neu erstellt;
eine ukrainische Originaldatei lag nicht vor. Die Gliederung und die englischen Ausgangsphrasen
folgen der bereits übertragenen spanischen Version. Alle 142 Phrasen und Ländernamen mit ihren
Referenzen sind enthalten; Nummernsprünge und der fehlende Abschnitt 7 bleiben erhalten.
Die vier sprachbezogenen englischen Einträge nennen jetzt „Ukrainian“ statt „Spanish“.

\medskip
\textbf{Sprache und Anrede.} Verwendet wird ukrainische Standardsprache in kyrillischer Schrift,
keine russische Übersetzung oder Umschrift. Die Anrede ist durchgehend höflich: „Ви“, „Вам“,
„Ваша“. Die Großschreibung kennzeichnet die persönliche höfliche Anrede.
„радий“ in 1.1 gilt für einen männlichen Sprecher; eine Sprecherin sagt „рада“.
„Я на пенсії“ in 2.11 ist unabhängig vom Geschlecht.
Die Formulierungen sind als praktische Sätze fürs Funkgespräch gedacht. Die Übersetzung und
Lautschrift wurden mit KI-Unterstützung erstellt; die Prüfung durch ukrainische Muttersprachlerinnen
und Muttersprachler, besonders mit Amateurfunkerfahrung, ist ausdrücklich erwünscht.

\medskip
\textbf{Bewusste Anpassungen.} „Добрий день“ ist sowohl die übliche Tagesbegrüßung als auch
am Nachmittag passend; eine wörtliche Nachmittagsbegrüßung wird nicht konstruiert.
Entsprechend wünscht 9.16 einen schönen Tag; „Добрий вечір“ und „гарного вечора“ beziehen sich auf den Abend.
Die starren Uhrzeitangaben der spanischen Vorlage entfallen. In 6.11 bedeutet „вже весь день“,
dass es bereits den ganzen Tag regnet. Die Mitgliedsnummer in 8.12 lautet einheitlich 1234,
wie in der englischen Ausgangsphrase; die spanische Übersetzung enthält dort abweichend 1224.
Die Frequenz wird als 3,558 МГц mit Dezimalkomma geschrieben; der Wert entspricht 3.558 MHz.
„село“, „містечко“ und „місто“ sind alternative Ortsbeschreibungen, keine drei gleichzeitigen
Behauptungen über Sittingbourne. Namen, Orte, Rufzeichen und Zahlen sind austauschbare Beispiele.

\medskip
\textbf{Fachbegriffe.} „позивний“ bezeichnet das Rufzeichen, „рапорт“ den Signalrapport,
„завади“ die Störungen und „діапазон“ das Frequenzband. „саморобна“ steht für Eigenbau.
„довгий дріт“ bezeichnet die Langdrahtantenne, „диполь“ den Dipol, „трикутна рамка“ die Delta-Loop,
„логоперіодична“ die logarithmisch-periodische und „спрямована“ die Richtantenne.
„Doublet“ bleibt in 3.43 als Dipol-Typbezeichnung erhalten. G5RV und Яґі bleiben Bauartnamen;
das ґ in Яґі steht für \ipa{ɡ}. „через бюро“ meint den Versand über das QSL-Büro,
„безпосередньо“ den direkten Versand. Der Abgleich mit ukrainischen Amateurfunktexten dient der
Terminologie; die neue Fassung ist keine von der Liga der Funkamateure der Ukraine bestätigte Übersetzung.

\medskip
\textbf{Kyrillisch und Lautschrift.} Die IPA-Spalte ist eine breite, manuell erstellte Aussprachehilfe,
keine enge Transkription einer Aufnahme. Die sechs Vokalqualitäten werden als \ipa{i, ɪ, ɛ, a, ɔ, u}
wiedergegeben. Unbetonte Vokale bleiben weitgehend in ihrer Grundqualität; regionale und
kontextabhängige Annäherungen sowie Satzmelodie werden nicht abgebildet.
\ipa{ʲ} bezeichnet Erweichung beziehungsweise Palatalisierung; vor і werden auch entsprechende
leichte Palatalisierungen von Lippen- und Hinterzungenlauten markiert. Ь ist kein eigener Vokal.
Г entspricht \ipa{ɦ}, ґ entspricht \ipa{ɡ}, х entspricht \ipa{x}, ж entspricht \ipa{ʒ},
ч entspricht \ipa{tʃ} und щ der Folge \ipa{ʃtʃ}. Ein Bindebogen in \ipa{t͡s, d͡z}
fasst die beiden Bestandteile einer Affrikate zusammen. В wird vereinfachend durchgehend als \ipa{ʋ}
notiert; seine kontextabhängigen w- beziehungsweise u-artigen Varianten werden nicht eigens markiert.
Doppelkonsonanten werden als doppelte Zeichen geschrieben, zum Beispiel \ipa{nnʲ}.
Die Transkription berücksichtigt einige gebräuchliche Angleichungen in Zahlwörtern,
aber nicht sämtliche Lautverbindungen über Wortgrenzen hinweg.

\medskip
\textbf{Eigennamen, Kürzel und Zahlen.} Пітер, Сіттінгборн und Дувр sind ukrainische
Namenswiedergaben; die Lautschrift folgt diesen Formen, einschließlich ihrer ukrainischen Endungen.
Icom, doublet und die lateinischen Rufzeichen und Kürzel werden hier mit englischer Aussprache
beziehungsweise englischen Buchstabennamen gelesen; die Ziffern darin werden ukrainisch gesprochen.
Das ist eine gewählte Lesekonvention, keine Vorgabe zum internationalen Buchstabieralphabet.
Hellschreiber ist deutsch wiedergegeben; Olivia hat eine ukrainisch angepasste Aussprache.
Die Webadressen behalten ihre originale Schreibweise; „крапка“ steht für den Punkt.
599 wird ziffernweise „п'ять дев'ять дев'ять“ gelesen; 3,558 „три кома п'ять п'ять вісім“;
1234 „тисяча двісті тридцять чотири“, 706 „сімсот шість“.
In 2.2 wird 1960 als Jahreszahl mit der passenden Form „шістдесятого року“ gesprochen.
Schrägstriche trennen Alternativen oder Listenelemente. \texttt{<\#\#\#\#\#>} ist ein Platzhalter
und erscheint in der Lautschrift als Auslassung.

\medskip
\textbf{Ländernamen.} Die historischen englischen Quellnamen bleiben zur Orientierung erhalten.
Ukrainisch werden gebräuchliche heutige Namen verwendet: „Чехія“ für Czech Republic,
„Північна Македонія“ für Macedonia, „Молдова“ für Moldavia und „Білорусь“ für Belorussia.
„Голландія“ steht für Holland, „Нідерланди“ für das gesamte Land Netherlands.
„Шпіцберген“ ist die gebräuchliche ukrainische Bezeichnung für den Archipel Svalbard.
Die Liste enthält Staaten, Gebiete und Landesteile; die Überschrift lautet daher „Країни та території“.

\medskip
\textbf{Quellen und Kompilierung.} Phonetische Grundlage: Bernd Pompino-Marschall, Elena Steriopolo
und Marzena Żygis, „Ukrainian“, Journal of the International Phonetic Association 47(3), 2017,
S. 349--357; DOI: 10.1017/S0025100316000372. Fachsprachlicher Abgleich: Liga der Funkamateure
 der Ukraine, uarl.org.ua, unter anderem „Multiband UKRAINE“ AWARD (MBUA).
Einzelne Wortbetonungen wurden mit dem Wörterbuch СЛОВНИК.ua (slovnyk.ua) abgeglichen.
UTF-8-Datei, zweimal mit XeLaTeX kompilieren; benötigt werden DejaVu Sans und die üblichen
LaTeX-Pakete der Präambel. Keine externen Bild- oder Datendateien erforderlich.
\end{document}
